\documentclass[aspectratio=169]{beamer}

\usetheme{Berlin}
%Outras opções de tema:
%AnnArbor
%Antibes
%Bergen
%Berkeley
%*Berlin
%Boadilla	
%CambridgeUS
%Copenhagen
%Darmstadt
%Goettingen
%PaloAlto
%Szeged
%Warsaw

\usecolortheme{metropolis}
%Outras opções de cores:
%*beaver
%beetle
%seahorse
%*spruce
%wolverine
\usepackage{appendixnumberbeamer}
\usepackage[numbers,sort&compress]{natbib}
\bibliographystyle{plainnat}
\usepackage{latexsym,amsmath,xcolor,multicol,booktabs,calligra}
\usepackage{graphicx,pstricks,listings,stackengine}
\usepackage[scale=2]{ccicons}
\usepackage[portuges]{babel}
\usepackage[utf8]{inputenc}
\usepackage{aeguill}
\usepackage{media9} %to insert GIF images
\usepackage{wrapfig}
\usepackage{hyperref}
\usepackage{multicol}

\usepackage{xspace}
\newcommand{\themename}{\textbf{\textsc{metropolis}}\xspace}
%Setagem pra cor dos hyperlinks
\hypersetup{
    colorlinks=true,
    linkcolor=[RGB]{0,128,128},
    linkbordercolor={blue},
    filecolor=magenta,
    urlcolor=blue,
    citecolor=.
}


\title{Segurança em Redes IoT}
\subtitle{Identificação, monitoramento e priorização de ameaças}
\date{\today}
\author{Antônio Marcos Souza Pereira, Samuel Paiva Bernardes}
\institute{DCC - UFJF}
% \titlegraphic{\hfill\includegraphics[height=1.5cm]{logo.pdf}}
\setbeamertemplate{caption}[numbered]

\begin{document}

\maketitle

\AtBeginSection[]
{
\begin{frame}
\frametitle{Sumário}
\tableofcontents[currentsection]
\end{frame}
}

\section{Introdução}
\begin{frame}{Motivação}
    \begin{itemize}
        \item Redes IoT colocam câmeras, sensores, fechaduras, controladores e medidores dentro da rede.
        \item Cada equipamento novo também vira um ponto que pode falhar, vazar tráfego ou ser usado por um invasor.
        \item Muitos desses dispositivos rodam por anos, com pouco recurso de hardware e atualização rara de firmware.
        \item Quando um dispositivo IoT para, o efeito pode sair da rede e chegar ao ambiente físico.
    \end{itemize}
\end{frame}

\begin{frame}{Problema de pesquisa}
    \begin{block}{Pergunta-problema}
        Como apoiar a tomada de decisão de administradores de redes IoT na identificação, monitoramento e priorização de ameaças de segurança?
    \end{block}

    \vspace{0.3cm}
    \begin{itemize}
        \item Detectar tráfego suspeito é só uma parte do trabalho.
        \item O administrador precisa decidir qual alerta vem primeiro.
        \item Em uma rede real, bloquear tudo pode derrubar sensores, serviços e rotinas que dependem deles.
    \end{itemize}
\end{frame}

\section{Contextualização}

\begin{frame}{O que torna IoT diferente?}
    \begin{itemize}
        \item A rede mistura fabricantes, protocolos e dispositivos com capacidades muito diferentes.
        \item Sensores e controladores podem ficar ligados por anos, inclusive em laboratórios, prédios e ambientes industriais.
        \item Muitos equipamentos não geram logs úteis e não aceitam agentes de segurança tradicionais.
        \item Uma falha pequena, repetida em muitos dispositivos, vira problema de rede.
    \end{itemize}
\end{frame}

\begin{frame}{Ameaças comuns}
    \begin{itemize}
        \item Ataques DoS e DDoS tentam tirar serviços do ar pelo volume de pacotes.
        \item Botnets associadas ao Mirai exploram senhas fracas e dispositivos expostos.
        \item Varreduras procuram portas abertas, serviços antigos e versões vulneráveis.
        \item Spoofing, força bruta e ataques web aparecem quando há interfaces de administração na rede.
    \end{itemize}
\end{frame}

\begin{frame}{Dataset de referência}
    \begin{itemize}
        \item O CIC-IoT-2023 reúne tráfego IoT capturado em ambiente de pesquisa.
        \item O conjunto traz tráfego benigno e ataques como DDoS, DoS, Mirai, Recon, Spoofing, Web e Brute Force.
        \item Ele ajuda a observar padrões de comunicação associados a comportamento normal e malicioso.
        \item Nesta apresentação, o dataset entra apenas para contextualizar o tema.
    \end{itemize}
\end{frame}

\section{Delimitação do Tema}

\begin{frame}{Da detecção à decisão}
    \begin{itemize}
        \item Administradores de rede recebem muitos eventos, alertas e indicadores ao mesmo tempo.
        \item Um alerta diz que algo pode estar errado.
        \item A prioridade mostra o que precisa de atenção antes.
        \item Em IoT, uma resposta apressada pode interromper sensores, automação ou serviços dependentes.
    \end{itemize}
\end{frame}

\begin{frame}{Recorte do trabalho}
    \begin{itemize}
        \item Tema: segurança de redes IoT.
        \item Objeto observado: tráfego de rede.
        \item Interesse prático: apoiar quem monitora, investiga e prioriza eventos.
    \end{itemize}
\end{frame}

\section{Conclusão}

\begin{frame}{Síntese do tema}
    \begin{itemize}
        \item Redes IoT colocam dispositivos simples em tarefas ocasionalmente complexas.
        \item O problema cresce quando muitos alertas chegam ao mesmo tempo.
        \item O tema deste trabalho é identificar, monitorar e priorizar ameaças nesse tipo de rede.
    \end{itemize}
\end{frame}

\begin{frame}{Referências}
    \footnotesize
    \begin{itemize}
        \item Canadian Institute for Cybersecurity. \textit{CIC IoT Dataset 2023}. \url{https://www.unb.ca/cic/datasets/iotdataset-2023.html}
        \item Neto et al. \textit{CICIoT2023: A Real-Time Dataset and Benchmark for Large-Scale Attacks in IoT Environment}. Sensors, 2023. \url{https://www.mdpi.com/1424-8220/23/13/5941}
    \end{itemize}
\end{frame}

\end{document}